\begin{thebibliography}{99}
\footnotesize\itemsep=0pt

\bibitem{Andrews-1974}
Andrews G.E., On the general {R}ogers--{R}amanujan theorem, \textit{Mem. Amer.
 Math. Soc.}, Vol.~152, American Mathematical Society, Providence, RI, 1974.

\bibitem{Andrews-1981}
Andrews G.E., Multiple {$q$}-series identities, \textit{Houston~J. Math.}
 \textbf{7} (1981), 11--22.

\bibitem{Berndt-III}
Berndt B.C., Ramanujan's notebooks. {P}art~{III},
 \href{https://doi.org/10.1007/978-1-4612-0965-2}{Springer}, New York, 1991.

\bibitem{Bressoud-1980}
Bressoud D.M., Analytic and combinatorial generalizations of the
 {R}ogers--{R}amanujan identities,
 \href{https://doi.org/10.1090/memo/0227}{\textit{Mem. Amer. Math. Soc.}}
 \textbf{24} (1980), 54~pages.

\bibitem{Cao-2023}
Cao Z., Wang L., Multi-sum {R}ogers--{R}amanujan type identities,
 \href{https://doi.org/10.1016/j.jmaa.2022.126960}{\textit{J.~Math. Anal.
 Appl.}} \textbf{522} (2023), 126960, 24~pages,
 \href{http://arxiv.org/abs/2205.12786}{arXiv:2205.12786}.

\bibitem{CKP09}
Cho B., Koo J.K., Park Y.K., Arithmetic of the
 {R}amanujan--{G}\"{o}llnitz--{G}ordon continued fraction,
 \href{https://doi.org/10.1016/j.jnt.2008.09.018}{\textit{J.~Number Theory}}
 \textbf{129} (2009), 922--947.

\bibitem{FrGa19}
Frye J., Garvan F., Automatic proof of theta-function identities, in Elliptic Integrals, Elliptic Functions and Modular Forms in Quantum Field Theory,
 \textit{Texts Monogr. Symbol. Comput.},
 \href{https://doi.org/10.1007/978-3-030-04480-0_10}{Springer}, Cham, 2019,
 195--258, \href{http://arxiv.org/abs/1807.08051}{arXiv:1807.08051}.

\bibitem{red}
Gasper G., Rahman M., Basic hypergeometric series, 2nd ed., \textit{Encyclopedia Math.
 Appl.}, Vol.~96,
 \href{https://doi.org/10.1017/CBO9780511526251}{Cambridge University Press},
 Cambridge, 2004.

\bibitem{Hardy-1937}
Hardy G.H., The {I}ndian {M}athematician {R}amanujan,
 \href{https://doi.org/10.2307/2301659}{\textit{Amer. Math. Monthly}}
 \textbf{44} (1937), 137--155.

\bibitem{Nahm-1994}
Nahm W., Conformal field theory, dilogarithms, and three-dimensional manifolds,
 \textit{Adv. Appl. Clifford Algebras} \textbf{4} (1994), suppl.~1, 179--191.

\bibitem{Nahm-1995}
Nahm W., Conformal field theory and the dilogarithm, in X{I}th {I}nternational
 {C}ongress of {M}athematical {P}hysics ({P}aris, 1994), International Press,
 Cambridge, MA, 1995, 662--667.

\bibitem{Nahm-2007}
Nahm W., Conformal field theory and torsion elements of the {B}loch group, in
 Frontiers in Number Theory, Physics, and Geometry.~{II},
 \href{https://doi.org/10.1007/978-3-540-30308-4_2}{Springer}, Berlin, 2007,
 67--132, \href{http://arxiv.org/abs/hep-th/0404120}{arXiv:hep-th/0404120}.

\bibitem{Ra77}
Rankin R.A., Modular forms and functions,
 \href{https://doi.org/10.1017/CBO9780511566035}{Cambridge University Press},
 Cambridge, 1977.

\bibitem{Ro94}
Robins S., Generalized {D}edekind {$\eta$}-products, in The {R}ademacher Legacy
 to Mathematics ({U}niversity {P}ark, {PA}, 1992), \textit{Contemp. Math.},
 Vol.~166, \href{https://doi.org/10.1090/conm/166/01645}{American Mathematical
 Society}, Providence, RI, 1994, 119--128.

\bibitem{Rogers-1894}
Rogers L.J., Second memoir on the expansion of certain infinite products, \href{https://doi.org/10.1112/plms/s1-25.1.318}{\textit{Proc.
 Lond. Math. Soc.}} \textbf{25} (1893/94), 318--343.

\bibitem{Schlosser-2023}
Schlosser M.J., Bilateral identities of the {R}ogers--{R}amanujan type,
 \href{https://doi.org/10.1090/btran/158}{\textit{Trans. Amer. Math. Soc.
 Ser.~B}} \textbf{10} (2023), 1119--1140,
 \href{http://arxiv.org/abs/1806.01153}{arXiv:1806.01153}.

\bibitem{Slater-1952}
Slater L.J., Further identities of the {R}ogers--{R}amanujan type,
 \href{https://doi.org/10.1112/plms/s2-54.2.147}{\textit{Proc. London Math.
 Soc.}} \textbf{54} (1952), 147--167.

\bibitem{Wang-2023}
Wang L., New proofs of some double sum {R}ogers--{R}amanujan type identities,
 \href{https://doi.org/10.1007/s11139-022-00654-5}{\textit{Ramanujan~J.}}
 \textbf{62} (2023), 251--272,
 \href{http://arxiv.org/abs/2203.15572}{arXiv:2203.15572}.

\bibitem{Wang-2024}
Wang L., Explicit forms and proofs of {Z}agier's rank three examples for
 {N}ahm's problem,
 \href{https://doi.org/10.1016/j.aim.2024.109743}{\textit{Adv. Math.}}
 \textbf{450} (2024), 109743, 46~pages,
 \href{http://arxiv.org/abs/2211.04375}{arXiv:2211.04375}.

\bibitem{Wang-2024-2}
Wang L., Identities on {Z}agier's rank two examples for {N}ahm's problem,
 \href{https://doi.org/10.1007/s40687-024-00460-z}{\textit{Res. Math. Sci.}}
 \textbf{11} (2024), 49, 34~pages,
 \href{http://arxiv.org/abs/2210.10748}{arXiv:2210.10748}.

\bibitem{Warnaar}
Warnaar S.O., Zudilin W., Dedekind's {$\eta$}-function and
 {R}ogers--{R}amanujan identities,
 \href{https://doi.org/10.1112/blms/bdr019}{\textit{Bull. Lond. Math. Soc.}}
 \textbf{44} (2012), 1--11,
 \href{http://arxiv.org/abs/1001.1571}{arXiv:1001.1571}.

\bibitem{Zagier-2007}
Zagier D., The dilogarithm function, in Frontiers in Number Theory, Physics, and Geometry.~{II},
 \href{https://doi.org/10.1007/978-3-540-30308-4_1}{Springer, Berlin}, 2007,
 3--65.

\end{thebibliography}
